\chapter{Introduction}
\label{ch:intro}

\section{Background}

AK\&VEN is a family brand of Korean-made socks. The name joins the first letters of my parents' names, \textbf{AK}ylbek
and \textbf{VEN}era, with an ampersand drawn as the Kyrgyz ram's-horn ornament (\textit{kochkor müyüz}) shaped as a heart. My family has traded at Dordoi Bazaar in Bishkek since the
early 2000s and has sold socks there since 2018, from a shipping container numbered 70-E. Dordoi is one of the largest
markets in Central Asia: thousands of containers stacked in long rows, where families shop for the season and traders
from the regions of Kyrgyzstan, from Kazakhstan, Uzbekistan and Russia buy goods for their own shops.

Three habits shape how socks are sold there, and every one of them is missing from an ordinary online shop:
\begin{itemize}
	\item \textbf{The price is agreed, not read.} A customer asks, the seller answers, and the price moves with the
	quantity, the season and the relationship. A fixed price tag can feel impersonal to these
	customers, and discount codes are unfamiliar to them.
	\item \textbf{People buy for the whole family.} One visit covers the father, the mother and the children, each in
	a different size; the larger the set, the lower the price per pair.
	\item \textbf{Shops buy by the case.} A case holds about 240 pairs in packs of 10 or 12, often mixed, and the price
	is given on request because it follows the dollar rate, the factory price and the delivery.
\end{itemize}

This thesis builds an integrated online shop and inventory system for the business, in which these habits can continue
online: a customer can ask for a better price in a conversation, plan one order for the whole family, mix any socks to
reach a lower price step, and collect the order at the container or have it delivered. Shops and companies can build a
case and send a request that the family answers. Behind the shop, the same system keeps the stock, the orders and the
sales of the business in one place.

\section{Related work}
\label{sec:related}

\paragraph{E-commerce platforms.} General platforms such as the open-source Shopizer~\cite{shopizer},
Medusa~\cite{medusa} and Saleor~\cite{saleor} offer catalogues, carts, orders and an administration interface. Their
architecture was a model for this project in two respects: a single API serves both the shop and the administration,
and the administration is only another client of that API with more rights. They assume fixed prices, however, with
discounts expressed as coupons or price lists; none of them treats the price as the result of a conversation, and they
are built for multi-tenant scale that a single family stall does not need.

\paragraph{Inventory management.} Business systems such as Odoo~\cite{odoo} model stock as a sequence of moves between
locations rather than as a single number, so that every quantity can be explained by the operations that produced it.
The same idea appears in software architecture as event sourcing~\cite{fowlereventsourcing}, in which every change is
stored as an event and the current state is derived from them. AK\&VEN adopts a lighter form suited to a single stall:
the current stock is kept as a number for fast reads, and every change is also written to a ledger with its reason in
the same database transaction. To prevent overselling, it relies on row-level locks in PostgreSQL~\cite{pglocking},
taken in a fixed order during checkout.

\paragraph{Negotiation dialogue.} Automated negotiation has been studied in natural-language processing for years.
Lewis et al.~\cite{lewis2017deal} trained end-to-end models to negotiate the division of items and found that they
learned to deceive; He et al.~\cite{he2018decoupling} separated the negotiation \emph{strategy} from the
\emph{generation} of the sentences, which made the agents' behaviour more controllable. This thesis follows the same
separation in an engineering setting: the strategy (how much may be given) is fixed by deterministic code and the
owners' settings, and the language model is used for the conversation and a proposal only.

\paragraph{Reliability of language-model applications.} Large language models follow instructions that arrive in their
input, including instructions hidden in user text. Greshake et al.~\cite{greshake2023} showed that applications which
let model output drive actions can be taken over by such prompt injection, and the OWASP Top 10 for LLM
applications~\cite{owaspllm} lists prompt injection, sensitive information disclosure and excessive agency among the
main risks. The common advice is to treat model output as untrusted input, to keep secrets out of the prompt and to
limit what the model can cause. AK\&VEN applies all three to money: the model never sees the limit or the cost price,
its proposal is validated by code, and it cannot state a price in its own words.

\paragraph{Grounding.} Retrieval-augmented generation~\cite{lewis2020rag} grounds model answers in retrieved documents
rather than in the model's general knowledge. The assistant here receives the facts of the product being discussed
(fabric, care, origin, pack size, stock, the public price ladder) selected for the question; the database already has a
vector table (pgvector~\cite{pgvector}) prepared for embedding-based retrieval in future work.

\paragraph{Progressive web applications.} A progressive web application~\cite{pwa} is a website that can be installed
like an app and can keep working offline with a service worker. For customers at a bazaar with an unreliable mobile
connection this matters more than a native app store presence, and one application can serve customers and staff.

\section{Motivation}

The motivation is practical and personal. Until now the business has been run by hand. Stock is counted at the stall
and kept in notebooks and in my parents' memory; orders arrive as messages on WhatsApp and Telegram; and how well a
model sells is known only from experience. This causes three recurring problems. First, the recorded stock and the
socks on the shelf drift apart, so a customer is sometimes promised a size that has already been sold. Second, orders
are fulfilled late because they are scattered across chats and depend on who is free to reply. Third, the family has
no clear view of which products, sizes and customers bring the most sales, which makes reordering from the Korean
suppliers a matter of guesswork.

An integrated system addresses all three: customers see the real assortment and stock, choose sizes from the shop's
own size chart and place orders in one place; the family sees orders, stock, incoming goods and requests in one list,
with every change of stock explained; and sales are summarised on a dashboard.

The difficulty is that bargaining is part of the brand, and a language model is the obvious tool for a conversation in
three languages, but it is not safe to let a model decide prices. A shop that can be talked into a 90\% discount loses
money; a shop that refuses to talk is not Dordoi. The thesis is about resolving that tension with a design in which
the model may talk freely while the price stays under the owners' control.

\section{Thesis objective}

The objective is to design, implement, test and evaluate a complete online shop for AK\&VEN in which:
\begin{itemize}
	\item customers can browse, choose sizes and colours, plan a family set, ask for a better price, pay and collect or
	receive their order, in English, Russian or Kyrgyz, on a phone, a tablet or a laptop;
	\item the owners and staff manage the catalogue, stock, prices, orders, requests and settings in an administration
	interface in the same three languages;
	\item stock stays correct under concurrent orders: every change is recorded with its reason, the last pair cannot
	be sold twice, reserved stock is released when it is not collected, and the owners are warned when a product runs
	low;
	\item no price below the owners' limit can ever reach an order, whatever the assistant says, and this guarantee is
	shown by automated tests and by an evaluation on simulated customers.
\end{itemize}

\section{Research questions}
\label{sec:rq}

\begin{enumerate}[\bfseries RQ1.]
	\item Can a language model hold a price negotiation for a real shop while a deterministic component guarantees that
	no discount above the owners' limit is ever applied?
	\item How does a language-model negotiator compare with a rule-based one on the same simulated customers, in terms of
	discounts given, sales made and how often the limit had to step in?
	\item Can the bazaar's habits (bargaining, family sets, price steps by quantity, wholesale by the case) be expressed
	in an online shop without making it harder to use than an ordinary one?
	\item Can the stock of a small shop stay consistent and explainable when online orders, reservations, cancellations,
	deliveries of new goods and manual corrections happen at the same time?
\end{enumerate}

\section{Scope}

Built in this thesis: the shop and the administration as one progressive web application; the REST API with
authentication and four roles; the catalogue with sections, cuts, colours, sizes, pictures and stock; inventory
management with a stock ledger, incoming goods, reservations and low-stock alerts; collection pricing; the negotiation
assistant with a rule-based and a language-model implementation; the bag and checkout with pickup at the stall, a
courier and cargo carriers, payment through a \emph{simulated} Apple Pay and Google Pay wallet or in person at
handover; guest checkout; email notifications; a sales dashboard and order export; family profiles, gift boxes,
reviews, saved products and wholesale requests; an audit log; the test suite; and deployment with Docker Compose and
continuous integration.

Not built, on purpose: a real card payment provider (the wallet moves no money), SMS notifications, taxes, a payment
rail through a Kyrgyz bank, demand forecasting and synchronisation with the family's messenger chats. These are
discussed as future work in Chapter~\ref{ch:sum}.

\section{Structure of the thesis}

Chapter~\ref{ch:spec} specifies the system: its users, use cases, and functional and non-functional requirements.
Chapter~\ref{ch:user} is the user documentation: how to install and start the shop, and how customers, staff, owners and
the super admin use it. Chapter~\ref{ch:dev} is the developer documentation: the architecture, the data model, the API,
security, the pricing and negotiation logic, and the design of the frontend. Chapter~\ref{ch:test} describes how the
system is tested and presents the evaluation of the two assistants. Chapter~\ref{ch:sum} answers the research questions
and lists future work.
